% Abstrak (Indonesia dan Inggris)
\babtp{ABSTRAK}
\begin{abstrak}
Jaringan Long Term Evolution (LTE) membagi sumber daya radio kepada pengguna melalui penjadwal paket, sehingga pemilihan algoritma penjadwalan memengaruhi throughput dan keadilan layanan. Penelitian ini bertujuan mengukur throughput total dan indeks fairness Jain pada algoritma Round Robin (RR), Proportional Fair (PF), dan Max C/I, serta menentukan algoritma dengan kompromi terbaik. Penelitian dilaksanakan dengan simulator ns-3 pada satu sel LTE berlebar pita 20 MHz dan jumlah pengguna 10 sampai 50 yang bergerak acak. Setiap skenario diulang sebanyak 10 kali dengan seed berbeda dan hasilnya dirata-ratakan. Hasil menunjukkan bahwa Max C/I menghasilkan throughput tertinggi, yaitu 40,7 Mbps pada 50 pengguna, diikuti PF (36,8 Mbps) dan RR (33,4 Mbps). Sebaliknya, indeks fairness tertinggi dicapai RR (0,90), diikuti PF (0,83), sedangkan Max C/I terendah (0,28). PF memberikan kompromi terbaik karena throughput-nya hanya 9,6 persen di bawah Max C/I dengan fairness yang jauh lebih tinggi. Disimpulkan bahwa PF paling sesuai bila diperlukan keseimbangan antara efisiensi dan keadilan.

\textbf{Kata kunci:} throughput, fairness, penjadwalan paket, LTE, ns-3.
\end{abstrak}

\begin{halamanbaru}\end{halamanbaru}
\babtp{ABSTRACT}
\begin{abstrak}
The Long Term Evolution (LTE) network allocates radio resources to users through a packet scheduler, so the choice of scheduling algorithm affects throughput and service fairness. This research aims to measure the total throughput and Jain's fairness index of the Round Robin (RR), Proportional Fair (PF), and Max C/I algorithms, and to determine the algorithm with the best compromise. The research was carried out using the ns-3 simulator on a single LTE cell with a 20 MHz bandwidth and 10 to 50 randomly moving users. Each scenario was repeated 10 times with different seeds and the results averaged. The results show that Max C/I yields the highest throughput, 40.7 Mbps at 50 users, followed by PF (36.8 Mbps) and RR (33.4 Mbps). In contrast, the highest fairness index was achieved by RR (0.90), followed by PF (0.83), while Max C/I was the lowest (0.28). PF provides the best compromise because its throughput is only 9.6 percent below Max C/I with a much higher fairness. It is concluded that PF is the most suitable when a balance between efficiency and fairness is required.

\textbf{Keywords:} throughput, fairness, packet scheduling, LTE, ns-3.
\end{abstrak}

\begin{halamanbaru}\end{halamanbaru}
